\documentclass[12pt]{article}
\usepackage[T1]{fontenc}
\usepackage[utf8]{inputenc}
\usepackage{lmodern}
\usepackage{textcomp}
\usepackage{enumitem}
\usepackage{geometry}
\usepackage{graphicx}
\usepackage{amsmath, amssymb}
\geometry{margin=1in}
\begin{document}
\begin{center}
Physics - Undergraduate Level Assessment 5 \
Total Marks: 40 \
Answer all questions.
\end{center}

\section*{Multiple Choice (10 marks)}
\begin{enumerate}[label=\arabic*.]
\item A grating spectrometer can just barely resolve two wavelengths of 500 nm and 502 nm, respectively. Which of the following gives the resolving power of the spectrometer?
\begin{enumerate}[label=(\alph*)]
    \item 2
    \item 250
    \item 5,000
    \item 10,000
\end{enumerate}
\item The emission spectrum of the doubly ionized lithium atom Li++ (Z = 3, A = 7) is identical to that of a hydrogen atom in which all the wavelengths are
\begin{enumerate}[label=(\alph*)]
    \item decreased by a factor of 9
    \item decreased by a factor of 49
    \item decreased by a factor of 81
    \item increased by a factor of 9
\end{enumerate}
\item De Broglie hypothesized that the linear momentum and wavelength of a free massive particle are related by which of the following constants?
\begin{enumerate}[label=(\alph*)]
    \item Planck's constant
    \item Boltzmann's constant
    \item The Rydberg constant
    \item The speed of light
\end{enumerate}
\item The best type of laser with which to do spectroscopy over a range of visible wavelengths is
\begin{enumerate}[label=(\alph*)]
    \item a dye laser
    \item a helium-neon laser
    \item an excimer laser
    \item a ruby laser
\end{enumerate}
\item Which of the following gives the total spin quantum number of the electrons in the ground state of neutral nitrogen (Z = 7)?
\begin{enumerate}[label=(\alph*)]
    \item 1/2
    \item 1
    \item 3/2
    \item 5/2
\end{enumerate}
\end{enumerate}

\section*{True / False (10 marks)}
\textit{Answer True or False}
\begin{enumerate}[label=\arabic*.]
\setcounter{enumi}{5}
\item True or False: The resolving power of the grating spectrometer for wavelengths of 500 nm and 502 nm is 250.
\item True or False: The magnetic susceptibility of a doped semiconductor can determine the sign of its charge carriers.
\item True or False: An electron in the n = 4, l = 1 state can transition to the n = 3, l = 2 state because it increases the angular momentum quantum number.
\item True or False: In a single-electron atom with l = 2, there are 5 allowed values for the magnetic quantum number m\_l: -2, -1, 0, 1, and 2.
\item True or False: In a reversible thermodynamic process, the entropy of the system and its environment remains unchanged.
\end{enumerate}

\section*{Long Form Response (20 marks)}
\textit{Answer each question with a clear and complete explanation.}
\begin{enumerate}[label=\arabic*.]
\setcounter{enumi}{10}
\item Calculate the total radiation force exerted on a completely reflecting surface of 3 $m^2$ at Earth's surface with a solar flux of 1,000 W/$m^2$. Discuss the underlying principles involved.
\item Discuss the effects of doubling the frequency of laser light on the separation of bright fringes observed in a double-slit experiment with a slit separation of 0.5 micrometers.
\end{enumerate}

\vfill
\noindent\textit{End of Paper}
\end{document}